\documentclass[12pt,letterpaper,reqno]{amsart}

\usepackage{epsfig}
\usepackage{amsmath}
\usepackage{amssymb}
\usepackage{amsthm}
\usepackage{indentfirst}
\usepackage{xspace}
\usepackage{multirow}
\usepackage{hyperref}
\usepackage{caption}
\usepackage{xcolor}
\usepackage{verbatim}
\usepackage[letterpaper,margin=1in,headheight=15pt]{geometry}
\usepackage{mathpazo}
\usepackage{tikz-cd}
\usepackage{booktabs}
\usepackage{framed}
\usepackage{float}
\usepackage{thmtools}
\usepackage{dashrule}
\usepackage{silence}
\WarningFilter{gitinfo2}{I can't find}
\WarningFilter{biblatex}{Patching}
\usepackage[missing=]{gitinfo2}
\usepackage{fancyhdr}
\usepackage{mathtools}
\usepackage{enumitem}

\setlist[enumerate]{topsep=0.4pt,itemsep=0.4ex,partopsep=0.4ex,parsep=0.4ex}

\mathtoolsset{showonlyrefs}

\setlength{\footskip}{18pt}

\definecolor{purple}{rgb}{0.4,0.05,0.8}
\definecolor{darkblue}{rgb}{0.1,0.1,0.6}
\definecolor{darkred}{rgb}{0.45,0.1,0.1}
\definecolor{darkgreen}{rgb}{0.0,0.3,0.06}
\hypersetup{colorlinks=true,urlcolor=darkred,linkcolor=purple,citecolor=darkred}
\definecolor{shadecolor}{rgb}{0.9,0.9,0.9}

% theorem environments
\declaretheoremstyle[spaceabove=0.25cm,spacebelow=0.25cm,notefont=\normalfont\bfseries, notebraces={(}{)}]{theorem}
\declaretheoremstyle[spaceabove=0.25cm,spacebelow=0.25cm,bodyfont=\normalfont,notefont=\normalfont\bfseries, notebraces={(}{)}]{noital}
\declaretheoremstyle[spaceabove=0.25cm,spacebelow=0.25cm,bodyfont=\normalfont\color{darkblue},notefont=\normalfont\bfseries, notebraces={(}{)}]{noitalblue}
\declaretheoremstyle[spaceabove=0.25cm,spacebelow=0.25cm,bodyfont=\normalfont\color{darkgreen},notefont=\normalfont\bfseries, notebraces={(}{)}]{green}
\declaretheoremstyle[spaceabove=0.25cm,spacebelow=0.25cm,bodyfont=\normalfont\color{darkred},notefont=\normalfont\bfseries, notebraces={(}{)}]{red}
\declaretheoremstyle[spaceabove=0.25cm,spacebelow=0.25cm,bodyfont=\normalfont,notefont=\normalfont\bfseries,qed=$\qedsymbol$,notebraces={(}{)}]{proofstyle}

\newcounter{pset}

\declaretheorem[name=Theorem,numberwithin=section,style=red]{thm}
\declaretheorem[name=Conjecture,numberwithin=section,style=theorem]{conj}
\declaretheorem[name=Proposition,sibling=thm,style=red]{prop}
\declaretheorem[name=Proposition,sibling=thm,numbered=no]{prop*}
\declaretheorem[name=Construction,sibling=thm,style=theorem]{constr}
\declaretheorem[name=Corollary,sibling=thm,style=red]{cor}
\declaretheorem[name=Lemma,sibling=thm,style=red]{lem}
\declaretheorem[name=Definition,sibling=thm,style=noitalblue]{defn}
\declaretheorem[name=Definition-Proposition,sibling=thm,style=noitalblue]{defprop}
\declaretheorem[name=Example,sibling=thm,style=green]{example}
\declaretheorem[name=Exercise,numberwithin=pset,style=noital]{exercise}
\declaretheorem[name=Solution,numberwithin=pset,style=noital]{solution}
% \declaretheorem[name=Proof,style=proofstyle,numbered=no]{pf}

\numberwithin{equation}{section}


% macros for convenience
\newcommand{\tops}{\texorpdfstring}

\newcommand{\nid}{\noindent}

\newcommand{\fa}{{\mathfrak a}}
\newcommand{\fp}{{\mathfrak p}}
\newcommand{\fk}{{\mathfrak k}}
\newcommand{\fg}{{\mathfrak g}}
\newcommand{\fh}{{\mathfrak h}}
\newcommand{\fn}{{\mathfrak n}}
\newcommand{\fq}{{\mathfrak q}}
\newcommand{\fm}{{\mathfrak m}}
\newcommand{\fr}{{\mathfrak r}}
\newcommand{\fu}{{\mathfrak u}}
\newcommand{\fG}{{\mathfrak G}}
\newcommand{\fgl}{{\mathfrak {gl}}}

\newcommand{\hC}{{\widehat C}}
\newcommand{\hS}{{\widehat S}}
\newcommand{\hX}{{\widehat X}}

\newcommand{\cC}{\ensuremath{\mathcal C}}
\newcommand{\cD}{\ensuremath{\mathcal D}}
\newcommand{\cW}{\ensuremath{\mathcal W}}
\newcommand{\cG}{\ensuremath{\mathcal G}}
\newcommand{\cB}{\ensuremath{\mathcal B}}
\newcommand{\cL}{\ensuremath{\mathcal L}}
\newcommand{\cS}{\ensuremath{\mathcal S}}
\newcommand{\cF}{\ensuremath{\mathcal F}}
\newcommand{\bcF}{\ensuremath{\overline{\mathcal F}}}
\newcommand{\cK}{\ensuremath{\mathcal K}}
\newcommand{\cZ}{\ensuremath{\mathcal Z}}
\newcommand{\cM}{\ensuremath{\mathcal M}}
\newcommand{\cP}{\ensuremath{\mathcal P}}
\newcommand{\cN}{\ensuremath{\mathcal N}}
\newcommand{\cO}{\ensuremath{\mathcal O}}
\newcommand{\cT}{\ensuremath{\mathcal T}}
\newcommand{\cH}{\ensuremath{\mathcal H}}
\newcommand{\cX}{\ensuremath{\mathcal X}}
\newcommand{\cY}{\ensuremath{\mathcal Y}}
\newcommand{\cA}{\ensuremath{\mathcal A}}
\newcommand{\cI}{\ensuremath{\mathcal I}}

\newcommand{\tphi}{{\ensuremath{\widetilde \phi}}}
\newcommand{\tf}{{\ensuremath{\widetilde f}}}
\newcommand{\tg}{{\ensuremath{\widetilde g}}}
\newcommand{\tiota}{{\ensuremath{\widetilde \iota}}}
\newcommand{\talpha}{{\ensuremath{\widetilde \alpha}}}
\newcommand{\tbeta}{{\ensuremath{\widetilde \beta}}}
\newcommand{\tU}{{\ensuremath{\widetilde U}}}
\newcommand{\tV}{{\ensuremath{\widetilde V}}}
\newcommand{\tM}{{\ensuremath{\widetilde M}}}
\newcommand{\tN}{{\ensuremath{\widetilde N}}}
\newcommand{\tomega}{{\ensuremath{\widetilde \omega}}}
\newcommand{\tF}{{\ensuremath{\widetilde F}}}

\newcommand{\R}{\ensuremath{\mathbb R}}
\newcommand{\C}{\ensuremath{\mathbb C}}
\newcommand{\PP}{\ensuremath{\mathbb P}}
\newcommand{\Z}{\ensuremath{\mathbb Z}}
\newcommand{\Q}{\ensuremath{\mathbb Q}}
\newcommand{\A}{\ensuremath{\mathbb A}}
\newcommand{\bbT}{\ensuremath{\mathbb T}}
\newcommand{\bbH}{\ensuremath{\mathbb H}}
\newcommand{\bbI}{\ensuremath{\mathbb I}}
\newcommand{\bbS}{\ensuremath{\mathbb S}}
\newcommand{\bbQ}{\ensuremath{\mathbb Q}}
\newcommand{\bbZ}{\ensuremath{\mathbb Z}}
\newcommand{\bbF}{\ensuremath{\mathbb F}}

\newcommand{\half}{\ensuremath{\frac{1}{2}}}
\newcommand{\qtr}{\ensuremath{\frac{1}{4}}}
\newcommand{\bq}{{\mathbf q}}
\newcommand{\bbN}{{\mathbb N}}
\newcommand{\N}{{\mathcal N}}
\newcommand{\F}{{\mathcal F}}
\newcommand{\HH}{{\mathcal H}}
\newcommand{\LL}{{\mathcal L}}
\newcommand{\RR}{{\mathcal R}}
\newcommand{\V}{{\mathcal V}}
\newcommand{\dirac}{\!\!\not\!\partial}
\newcommand{\Dirac}{\!\!\not\!\!D}
\newcommand{\cE}{{\mathcal E}}
\newcommand{\vs}{\not\!v}
\newcommand{\kahler}{K\"ahler\xspace}
\newcommand{\painleve}{Painlev\'e\xspace}
\newcommand{\kq}{/\!\!/}
\newcommand{\kql}[1]{/\!\!/\!\!_#1\,}
\newcommand{\hk}{hyperk\"ahler\xspace}
\newcommand{\Hk}{Hyperk\"ahler\xspace}
\newcommand{\hkq}{/\!\!/\!\!/\!\!/}
\newcommand{\hkql}[1]{/\!\!/\!\!/\!\!/\!\!_#1\,}
\newcommand{\del}{\ensuremath{\partial}}
\newcommand{\delbar}{\ensuremath{\overline{\partial}}}
\newcommand{\vphi}{{\vec \phi}}
\newcommand{\I}{{\mathrm i}}
\newcommand{\J}{{\mathrm j}}
\newcommand{\K}{{\mathrm k}}
\newcommand{\e}{{\mathrm e}}
\newcommand\bid{{\mathbf 1}}
\newcommand{\De}{\mathrm{D}}
\newcommand{\de}{\mathrm{d}}
\newcommand{\ab}{\mathrm{ab}}
\newcommand{\vol}{\mathrm{vol}}
\renewcommand{\sf}{\mathrm{sf}}
\newcommand{\inst}{\mathrm{inst}}
\newcommand{\eff}{\mathrm{eff}}
\newcommand{\fiber}{\mathrm{fiber}}
\newcommand{\rmvert}{\mathrm{vert}}
\newcommand{\rmcyc}{\mathrm{cyc}}
\newcommand{\closed}{\mathrm{closed}}
\newcommand{\exact}{\mathrm{exact}}
\newcommand{\op}{\mathrm{op}}
\newcommand{\p}{\partial}
\newcommand{\GL}{\mathrm{GL}}
\newcommand{\SL}{\mathrm{SL}}
\newcommand{\Mat}{\mathrm{Mat}}
\newcommand{\PGL}{\mathrm{PGL}}
\newcommand{\PSL}{\mathrm{PSL}}
\newcommand{\formal}{\mathrm{formal}}

\newcommand{\abs}[1]{\lvert#1\rvert}
\newcommand{\norm}[1]{\lVert#1\rVert}
\newcommand{\IP}[1]{\langle#1\rangle}
\newcommand{\DIP}[1]{\langle\!\langle#1\rangle\!\rangle}
\newcommand{\dwrt}[1]{\frac{\partial}{\partial#1}}
\newcommand{\eps}{\epsilon}
\newcommand{\simarrow}{\xrightarrow\sim}

\newcommand{\as}{{\text{ as }}}
\newcommand{\tin}{{\text{ in }}}

\newcommand{\Airy}{\textrm{Airy}}
\newcommand{\Betti}{\textrm{Betti}}
\newcommand{\Conf}{\textrm{Conf}}
\newcommand{\Fl}{\textrm{Fl}}
\newcommand{\Sol}{\textrm{Sol}}
\newcommand{\Mon}{\textrm{Mon}}
\newcommand{\dR}{\textrm{dR}}
\newcommand{\dRp}{\textrm{dR,p}}
\newcommand{\Ai}{\textrm{Ai}}
\newcommand{\Bi}{\textrm{Bi}}
\newcommand{\tC}{{\widetilde C}}
\newcommand{\tSigma}{{\widetilde \Sigma}}
\newcommand{\tS}{{\widetilde S}}
\newcommand{\near}{\textrm{near}}
\newcommand{\far}{\textrm{far}}
\newcommand{\std}{\textrm{std}}

\newcommand{\pd}[2]{\frac{\partial #1}{\partial #2}}

\newcommand{\mmaref}[1]{}

\newcommand{\ti}[1]{\textit{#1}}
\newcommand{\tb}[1]{\textbf{#1}}

\DeclareMathOperator{\sgn}{sgn}
\DeclareMathOperator{\supp}{supp}
\DeclareMathOperator{\ad}{ad}
\DeclareMathOperator{\im}{Im}
\DeclareMathOperator{\re}{Re}
\DeclareMathOperator{\Tr}{Tr}
\DeclareMathOperator{\Gr}{Gr}
\DeclareMathOperator{\Int}{Int}
\DeclareMathOperator{\Bd}{Bd}
\DeclareMathOperator{\End}{End}
\DeclareMathOperator{\Hom}{Hom}
\DeclareMathOperator{\Aut}{Aut}
\DeclareMathOperator{\Sym}{Sym}
\DeclareMathOperator{\Lie}{Lie}
\DeclareMathOperator{\diag}{diag}
\DeclareMathOperator{\Bun}{Bun}
\DeclareMathOperator{\Vect}{Vect}
\DeclareMathOperator{\Span}{Span}
\DeclareMathOperator{\grad}{grad}
\DeclareMathOperator{\rank}{rank}
\DeclareMathOperator{\ind}{ind}
\DeclareMathOperator{\coker}{coker}
\DeclareMathOperator{\Jac}{Jac}
\DeclareMathOperator{\Hol}{Hol}
\DeclareMathOperator{\Res}{Res}
\DeclareMathOperator{\gr}{gr}
\DeclareMathOperator{\Ob}{Ob}
\DeclareMathOperator{\Path}{Path}
\DeclareMathOperator{\LPath}{LPath}
\DeclareMathOperator{\id}{id}

\newcommand{\insfig}[3]{

\medskip
\noindent
\begin{minipage}{\linewidth}

\captionsetup{type=figure} 

\makebox[\linewidth]{\includegraphics[keepaspectratio=true,scale=#2]{figures/#1.pdf}}

\captionof{figure}{#3}
\label{fig:#1}

\end{minipage}
\medskip

}


% \newcommand{\insfig}[2]{\begin{figure}[htbp] \centering \includegraphics[scale=#2]{figures/#1-crop.pdf} \label{fig:#1} \end{figure}}
% syntax: \insfig{name}{0.5}{caption}

\newcommand{\warn}[1]{{\color{red}{[#1]}}}
\newcommand{\fixme}[1]{{\color{orange}{[#1]}}}
\newcommand{\currentposition}{{\color{blue} \noindent\makebox[\linewidth]{\hdashrule{\paperwidth}{1pt}{3mm}}}}

% \mathtoolsset{showonlyrefs}

\begin{document}

\makeatletter
\renewcommand{\sectionautorefname}{\S\@gobble}
\renewcommand{\subsectionautorefname}{\S\@gobble}
\renewcommand{\subsubsectionautorefname}{\S\@gobble}
\makeatother

\pagestyle{fancy}
\lhead{{\tiny \color{gray} \tt \gitAuthorIsoDate}}
\chead{\tiny \ti{MATH 2550: Analysis 1, Fall 2026, Yale University}}
\rhead{{\tiny \color{gray} \tt \gitAbbrevHash}}
\cfoot{}
\renewcommand{\headrulewidth}{0.5pt}

\thispagestyle{fancy}

\newpage

\setcounter{pset}{1}
\setcounter{exercise}{0}

\begin{center}
{\bf{Problem set \arabic{pset}}} \\
{\bf{Due Thursday January 22, 11:59pm}}
\end{center}


Problem sets will be assigned weekly. You should submit your solutions to the problem sets 
in PDF form via the Gradescope interface in Canvas.
LaTeX is a very nice tool to use for this purpose, but handwritten
solutions are good too as long as they are clearly legible.

On the problem sets, you are free --- and indeed strongly encouraged --- to collaborate with other
students in our class.
Discussing mathematics with others is the easiest 
way of learning it.

There are two rules:

\begin{enumerate}

\item \textit{You are not permitted to post problems from the problem sets to publicly accessible websites.}

\item \textit{You must write down your solutions to the problems on your own: thus 
your solutions should not be identical or near-identical to solutions submitted by other students, and they also should not look like output from ChatGPT or similar tools.}
\end{enumerate}

You are of course free and encouraged to discuss the problems with me or with our TA or ULAs.
You are also free to use other resources, such as our textbook, our lecture notes,
any other books, information found on the Internet, or AI tools.
(Despite this last point, I do not suggest searching
the Internet for solutions to the problems, or using AI to solve the problems; 
even if you do get correct solutions, you won't learn much that way.)

In writing up your solutions, you are free to use without proof anything stated in class
or in the lecture notes, or anything in the course text. 
You should be sure to state clearly what you are using.
If you have some doubt about whether something is OK to use, just ask!

\smallskip

% \noindent\hrulefill


\begin{exercise}[5 points] Suppose $A$, $B$, and $C$ are sets, $f:A\to B$ is surjective, and $g:B\to C$ is surjective. Show that $g\circ f$ is surjective. 
\end{exercise}


\begin{exercise}[10 points; Rudin 1.3] \label{ex:field-checks}
Suppose $F$ is a field with $x,y,z \in F$.
Prove carefully from the field axioms:
\begin{enumerate}
  \item If $x \neq 0$ and $xy = xz$, then $y = z$.
  \item If $x \neq 0$ and $xy = x$, then $y = 1$.
  \item If $xy = 1$, then $x \neq 0$ and $y = x^{-1}$.
  \item If $x \neq 0$, then $(x^{-1})^{-1} = x$.
\end{enumerate}
% (For the purposes of this exercise, you should use only the properties given in the definition of a field. For the rest of the course, though, we will allow ourselves freely to use the usual rules of arithmetic, without carefully justifying them from the definition of a field every time.)
\end{exercise}


%\begin{exercise}[5 points; {Rudin 1.1}]
%For this exercise, all you need to know about
%$\R$ is that it is a field which contains the field $\Q$.
%
%Suppose $r \in \Q$, $r \neq 0$, and $x \in \R$, $x \notin \Q$.  Prove that $r+x \notin \Q$ 
%and $rx \notin \Q$.
%\end{exercise}


\begin{exercise}[10 points] In this problem you may freely use without proof all the standard facts about 
arithmetic of rational numbers and integers.

\begin{enumerate}

\item Prove that $\bbQ$ does not contain any $x$ with $x^2 = 3$. (For this part it may be useful to use the
fundamental theorem of arithmetic, which you may use without proof.)

\item Let $\bbQ(\sqrt{3})$ be the set of ordered pairs
$(a,b)$ with $a, b \in \bbQ$.
Informally, we think of each pair $(a,b)$
as representing a number ``$a + b \sqrt{3}$.''
We can equip $\bbQ(\sqrt{3})$ with addition and product laws as follows:
% \begin{gather}
%   (a + b \sqrt{3})  +  (a' + b' \sqrt{3}) = (a+a') + (b+b')\sqrt{3} \, , \\
%   (a + b \sqrt{3}) (a' + b' \sqrt{3}) = (aa' + 3bb') + (ab' + ba') \sqrt{3} \, .
% \end{gather}
\begin{gather}
  (a,b)  +  (a',b') = (a+a', b+b') \, , \\
  (a,b) \cdot (a',b') = (aa' + 3bb', ab' + ba') \, .
\end{gather}
(These rules are motivated by applying the usual rules of arithmetic to sums of the form ``$a + b \sqrt3$''.)

Show that $\bbQ(\sqrt{3})$ can be made into a field, with these addition and product laws.
(This means saying carefully what are the $0$ and $1$ elements, what is the negation law, 
what is the inversion law, and 
then proving that all the axioms of a field are satisfied.
This will probably be the lengthiest part of your writeup.)

\item In solving part (2), did you use the result of part (1)? If so, where?

\item $\bbQ(\sqrt3)$ contains an element $3$, defined as $3 = 1 + 1 + 1$, where $1$ is the $1$ element
you already defined in part (2). What is the element $3$? Prove that $\bbQ(\sqrt3)$ contains exactly two elements $x$ with $x^2 = 3$.

\end{enumerate}

\end{exercise}


% \begin{exercise}[5 points]
% As we have stated, the usual order relation on $\bbQ$ can be expressed as:
% \begin{equation}
%   p/q < p'/q' \quad \iff \quad pq' < p'q \, 
% \end{equation}
% where $q$ and $q'$ are chosen to be positive (note that this can always be done). 

% We can define a different relation $\prec$, as follows:
% assume that all fractions are written in lowest terms and with positive denominators. Define
% \begin{equation}
%   p/q \prec p'/q' \quad \iff \quad pq < p'q' \, .
% \end{equation}
% Prove that the relation $\prec$ does \ti{not} make $\bbQ$ into an ordered set.

% \end{exercise}

\begin{exercise}[8 points]
 
Let $A$ be a nonempty subset of an ordered set $S$. Suppose $\alpha$ is a lower bound of $A$ and
$\beta$ is an upper bound of $A$. 

\begin{enumerate}
\item Prove that $\alpha \le \beta$.
\item Note that in the statement of the problem we require $A$ to be nonempty. Would the statement $\alpha \le \beta$ still hold even
if we allow $A$ to be empty? Explain why or why not. In solving part (1), did you use the fact that $A$ is nonempty?
If so, where did you use it?
\end{enumerate}

\end{exercise}


%\begin{exercise}
%Prove that an ordered field cannot have a maximum element.
%\end{exercise}

% \begin{exercise}[{Rudin 1.5}]
% Let $A \subset \R$ be nonempty and bounded below.  Let $-A = \{-x\ \vert\ x \in A\}$.
% Prove that $\inf(A) = - \sup(-A)$.
% \end{exercise}

% \begin{exercise}
% Does the empty set $\emptyset \subset \Q$ have a supremum in $\Q$? Prove your answer.
% \end{exercise}

% \begin{exercise}
% Suppose $A \subset \R$ is nonempty.
% \begin{enumerate}
% \item Suppose $L = \sup(A)$. Then, for all $\eps \in \R$ with $\eps > 0$, 
% prove that there exists some $x \in A$ such that $L - \eps < x \le L$.
% \item Suppose $L$ is an upper bound for $A$ and $L \in A$.
% Prove that $L = \sup(A)$.
% \end{enumerate}
% \end{exercise}

% \begin{exercise}
% Suppose $a, b \in \R$, with $a < b$.  Let $A = (a,b)$.
% Prove that $\sup (A) = b$.  (Similarly $\inf (A) = a$, but you don't need to write
% out the proof of this part.)
% \end{exercise}


% \begin{exercise}
% This exercise is a reminder that the ``size'' of infinite sets can
% be a bit counterintuitive.

% \begin{enumerate}
%   \item
% Suppose $a, b \in \R$, with $a < b$.
% Prove that there is a bijection from $(a,b)$ to $(0,1)$.

% \item Prove that there is a bijection from $(0,1)$ to $(1,\infty)$.

% \item Suppose $a,b \in \R$. 
% Prove that there is a bijection from $(a,\infty)$ to $(b,\infty)$.
% \end{enumerate}
% \end{exercise}

\end{document}

